\chapter{Python Primer} \label{python_primer}

If you are new to programming or need a refresher, this primer covers the essential Python
and MicroPython concepts used throughout this workshop.

The primer is split into two parts:
\begin{enumerate}
    \item \textbf{General Python Fundamentals:} Core syntax, variables, functions, loops, and logic.
    \item \textbf{MicroPython on Microcontrollers:} How Python controls hardware pins, reads sensors,
    runs continuous event loops, and uses helper classes.
\end{enumerate}

\begin{tcolorbox}[colback=blue!5!white,colframe=blue!40!black]
    \textbf{MicroPython vs. Standard Python (CPython):}
    \href{https://micropython.org/}{MicroPython} is a lightweight, efficient implementation of
    Python 3 optimized to run directly on microcontrollers like the ESP32-C3. It includes built-in
    modules like \texttt{machine} for controlling hardware (pins, PWM, analog inputs, and I\textsuperscript{2}C).

    While some large desktop libraries (such as \texttt{pandas} or \texttt{pygame}) are not available
    in MicroPython, core Python syntax, data structures, and control flow are identical. What you learn
    here applies directly to general Python programming!
\end{tcolorbox}


\section{Part 1: General Python Fundamentals}

\subsection{A simple Python script}
Here is a small Python example demonstrating variables, functions, loops, and conditions:

\begin{lstlisting}[language=Python,caption=A basic Python script]
def show_message(text, repeat=1):
    """Prints a message to the console with iteration numbers."""
    if repeat > 1:
        print("Repeating", repeat, "times:")
    
    for i in range(repeat):
        print("Iteration", i, ":", text)

# Main program flow
user_name = "Emily"
show_message(user_name)
show_message(user_name, repeat=3)
\end{lstlisting}

\subsection{Breaking down the syntax}

\begin{itemize}
    \item \textbf{Defining functions (\texttt{def}):} On line 1, \texttt{def show\_message(text, repeat=1):}
    defines a reusable function. The parameter \texttt{text} is required, while \texttt{repeat=1} is an
    \textbf{optional parameter} with a default value of 1.
    
    \item \textbf{Docstrings and comments:} Lines 2 is a docstring (\texttt{"""..."""}) describing what the
    function does. Lines starting with \texttt{\#} (like line 9) are comments ignored by Python, written
    to help humans understand the code.
    
    \item \textbf{Indentation matters:} Python does not use curly brackets (\texttt{\{ \}}) to enclose code blocks.
    Instead, it uses \textbf{indentation} (typically 4 spaces). All lines indented under \texttt{def},
    \texttt{if}, or \texttt{for} belong to that block.
    
    \item \textbf{Conditionals (\texttt{if}):} Line 3 checks if \texttt{repeat > 1}. If true, it executes the indented
    \texttt{print} statement.
    
    \item \textbf{Loops and \texttt{range()}:} Line 6 starts a \texttt{for} loop. In Python, \texttt{range(3)}
    generates the sequence of numbers \texttt{0, 1, 2} (starting at 0 and running for 3 steps). The variable
    \texttt{i} takes on each value in turn.
    
    \item \textbf{Variables and function calls:} Line 10 assigns the string \texttt{"Emily"} to the variable
    \texttt{user\_name}. Line 11 calls \texttt{show\_message} with default repeat (1 time). Line 12 calls it
    with \texttt{repeat=3}, which prints 3 iterations.
\end{itemize}

Running this script produces the following console output:
\begin{verbatim}
Iteration 0 : Emily
Repeating 3 times:
Iteration 0 : Emily
Iteration 1 : Emily
Iteration 2 : Emily
\end{verbatim}


\section{Part 2: MicroPython on Hardware}

\subsection{The microcontroller program anatomy}
Unlike a desktop script that runs once and exits, a microcontroller typically runs a continuous
\textbf{event loop} to monitor sensors and update outputs.

Here is an excerpt illustrating the standard hardware pattern from Project 1:

\begin{lstlisting}[language=Python,caption=A typical MicroPython hardware script]
import time
from machine import Pin, ADC, PWM

# Constants
ADC_MAX = 4095
PWM_MAX = 1023

# Hardware configuration
sensor = ADC(Pin(2))
sensor.atten(ADC.ATTN_11DB)
led = PWM(Pin(20), freq=1000)

try:
    while True:
        # Read the light sensor (0 to 4095)
        level = sensor.read()

        # Scale sensor reading to LED PWM duty (0 to 1023)
        duty = (ADC_MAX - level) * PWM_MAX // ADC_MAX
        led.duty(duty)

        print("Sensor:", level, "PWM Duty:", duty)
        time.sleep(0.05)

except KeyboardInterrupt:
    # Clean up when the stop button is pressed in the IDE
    led.duty(0)
    led.deinit()
\end{lstlisting}

\subsection{Key hardware programming patterns}

\begin{itemize}
    \item \textbf{The \texttt{machine} module:} Provides drivers for physical pins. \texttt{Pin(20)} creates a GPIO
    controller; wrapping it in \texttt{PWM(...)} enables pulse-width modulation for brightness or sound;
    wrapping it in \texttt{ADC(...)} enables analog voltage measurement.
    
    \item \textbf{The continuous \texttt{while True:} loop:} Microcontrollers do not have operating systems with
    windows and mice. The program enters a loop that repeats indefinitely, reading inputs and controlling outputs.
    
    \item \textbf{Integer division (\texttt{//}):} MicroPython uses \texttt{//} for integer division, which discards
    the remainder. This is useful for calculating whole-number PWM duty values or pixel positions.
    
    \item \textbf{Pacing with \texttt{time.sleep()}:} A brief pause (such as 50 milliseconds) prevents the loop from
    consuming 100\% of the CPU and keeps the console output readable.
    
    \item \textbf{Clean exits with \texttt{try / except KeyboardInterrupt}:} When you press the red Stop button in
    the Viper IDE, it sends a \texttt{KeyboardInterrupt} signal. Wrapping the main loop in \texttt{try / except}
    allows the program to turn off LEDs and release hardware resources cleanly.
\end{itemize}

\subsection{Object-Oriented Programming (Classes) in MicroPython}
As your projects grow, grouping data (state) and functions (behavior) into \textbf{classes} keeps your code
organized. For example, the \texttt{DebouncedButton} class used in Projects 2, 3, and 4 tracks pin state and
debounce timers:

\begin{lstlisting}[language=Python,caption=A simple helper class for a pushbutton]
class SimpleButton:
    def __init__(self, pin_num, action):
        """Initializes the pin with a pull-up resistor and stores a callback."""
        self.pin = Pin(pin_num, mode=Pin.IN, pull=Pin.PULL_UP)
        self.action = action
        self.press_count = 0

    def check(self):
        """Checks if the button was pressed (LOW when using pull-up)."""
        if self.pin.value() == 0:
            self.press_count += 1
            self.action(self.press_count)
\end{lstlisting}

\begin{itemize}
    \item \textbf{\texttt{\_\_init\_\_} method:} Automatically called when a new object is created (e.g.,
    \texttt{btn = SimpleButton(21, on\_press)}). It initializes variables unique to that instance.
    \item \textbf{\texttt{self}:} Refers to the specific object instance, allowing methods to access and modify
    its own properties (like \texttt{self.press\_count}).
    \item \textbf{Callbacks (functions as arguments):} Passing a function to \texttt{action} lets the button
    trigger your custom logic whenever a press is detected.
\end{itemize}


\section{MicroPython Concepts Reference Map}

Here is a quick reference for Python features you will encounter in each project:

\begin{table}[H]
\centering
\begin{tabular}{|l|p{6.5cm}|l|}
    \hline
    \textbf{Concept} & \textbf{What it does} & \textbf{Where used} \\
    \hline
    \textbf{Tuples \& Lists} & Ordered collections like RGB colors \texttt{(255, 0, 0)} and indexable lists \texttt{pixels[0]} & Projects 2, 4, 5 \\
    \hline
    \textbf{Callbacks} & Passing a function name as an argument to run on button events & Projects 2, 3, 4 \\
    \hline
    \textbf{Millisecond Timers} & Measuring reaction times and delays using \texttt{time.ticks\_ms()} and \texttt{time.ticks\_diff()} & Projects 2, 5 \\
    \hline
    \textbf{State Variables} & Using variables to track game phases (\texttt{STATE\_RED}, \texttt{STATE\_GREEN}) & Project 2 \\
    \hline
    \textbf{Hardware Interrupts} & Rapidly detecting rotary encoder clicks without polling & Project 3 \\
    \hline
    \textbf{Dictionaries \& JSON} & Key-value data mapping (\texttt{\{"mode": "rainbow"\}}) for network APIs & Project 5 \\
    \hline
    \textbf{I\textsuperscript{2}C Communication} & Communicating with multiple digital sensors over a shared 2-wire bus & Project 7 \\
    \hline
\end{tabular}
\caption{Python and MicroPython concepts used across the workshop projects.}
\label{tab:python_concepts}
\end{table}
